\documentclass[11pt,a4paper]{article}
\usepackage[utf8]{inputenc}
\usepackage[vietnamese]{babel}
\usepackage{amsmath,amssymb}
\usepackage{geometry}
\geometry{a4paper, total={180mm,260mm}, left=15mm, right=15mm, top=18mm, bottom=18mm}
\usepackage{tikz}
\usetikzlibrary{shapes.geometric, arrows.meta, positioning, calc}
\usepackage{tcolorbox}
\tcbuselibrary{skins,breakable}
\usepackage{xcolor}
\usepackage{multicol}
\usepackage{array}
\usepackage{fontawesome}
\usepackage{fancyhdr}

\pagestyle{fancy}
\fancyhf{}
\renewcommand{\headrulewidth}{0.5pt}
\renewcommand{\footrulewidth}{0.5pt}
\fancyhead[L]{\small\bfseries\color{teal}\faBook\ SỔ TAY KIẾN THỨC TOÁN 3}
\fancyhead[R]{\small\bfseries\color{teal}Lớp Toán Thầy Lĩnh -- loptoanthaylinh.vn}
\fancyfoot[C]{\small\color{gray}-- Sổ tay kiến thức: Bài 34. Thực hành trải nghiệm đo --}

\definecolor{primary}{HTML}{005F73}
\definecolor{secondary}{HTML}{0A9396}
\definecolor{accent}{HTML}{EE9B00}
\definecolor{danger}{HTML}{AE2012}
\definecolor{boxbg}{HTML}{F4F6F8}

\begin{document}

% --- BANNER TIÊU ĐỀ ---
\begin{tcolorbox}[
  enhanced,
  colback=white,
  colframe=primary,
  arc=3mm,
  boxrule=1.2pt,
  left=4mm, right=4mm, top=2.5mm, bottom=2.5mm,
  drop shadow=black!10!white
]
  \centering
  {\Large\bfseries\color{primary}\faBookmark\ BÀI 34: THỰC HÀNH TRẢI NGHIỆM ĐO}\\[1.5mm]
  {\small\itshape\color{black!75}Kĩ năng chọn đơn vị đo, ước lượng và thực hành đo độ dài, khối lượng, dung tích, nhiệt độ trên 1 trang A4}
\end{tcolorbox}

\vspace{-2mm}


\noindent
\begin{minipage}[t]{0.49\textwidth}\vspace{0pt}
\begin{tcolorbox}[
  enhanced,
  colback=blue!4!white,
  colframe=blue!75!black,
  arc=2mm, boxrule=0.8pt,
  left=2mm, right=2mm, top=1.8mm, bottom=1.8mm,
  title={\footnotesize\bfseries\color{white}\faKey\ 1. Nguyên tắc chọn đơn vị đo hợp lý}
]
{\footnotesize
\begin{itemize}\setlength{\itemsep}{1.5pt}\setlength{\leftskip}{-2.5mm}
  \item \textbf{Độ dài:}
  \begin{itemize}\setlength{\itemsep}{0.5pt}\setlength{\leftskip}{-2.5mm}
    \item Siêu nhỏ (bề dày sách, ngòi bút): dùng $\mathbf{mm}$.
    \item Trung bình (bút chì, quyển vở): dùng $\mathbf{cm, dm}$.
    \item Lớn (phòng học, sân trường): dùng $\mathbf{m}$.
    \item Khoảng cách giữa các tỉnh: dùng $\mathbf{km}$.
  \end{itemize}
  \item \textbf{Khối lượng:} Hạt gạo, quả trứng, gói bánh: $\mathbf{g}$; Bao gạo, con bò, người lớn: $\mathbf{kg}$.
  \item \textbf{Dung tích:} Thìa thuốc, hộp sữa chua: $\mathbf{ml}$; Thùng sơn, bình nước to: $\mathbf{l}$.
\end{itemize}
}
\end{tcolorbox}
\end{minipage}
\hfill
\begin{minipage}[t]{0.49\textwidth}\vspace{0pt}
\begin{tcolorbox}[
  enhanced,
  colback=teal!4!white,
  colframe=teal!75!black,
  arc=2mm, boxrule=0.8pt,
  left=2mm, right=2mm, top=1.8mm, bottom=1.8mm,
  title={\footnotesize\bfseries\color{white}\faSliders\ 2. Kĩ năng ước lượng thực tế}
]
{\footnotesize
\begin{itemize}\setlength{\itemsep}{1.5pt}\setlength{\leftskip}{-2.5mm}
  \item \textbf{Ước lượng:} Nhìn bằng mắt thường và dựa vào các vật mẫu quen thuộc để dự đoán số đo trước khi đo thực tế.
  \item \textbf{Thước đo vật mẫu quen thuộc:}
  \begin{itemize}\setlength{\itemsep}{0.5pt}\setlength{\leftskip}{-2.5mm}
    \item Bề dày đồng xu: khoảng $1\text{ mm} - 2\text{ mm}$.
    \item Chiều dài gang tay học sinh: khoảng $15\text{ cm}$.
    \item Một bước chân học sinh: khoảng $40\text{ cm} - 50\text{ cm}$.
    \item Một quả trứng gà: nặng khoảng $50\text{ g}$.
  \end{itemize}
\end{itemize}
}
\end{tcolorbox}
\end{minipage}

\vspace{-2.5mm}
% --- KHỐI 3: SƠ ĐỒ TƯ DUY TIKZ ---
\noindent
\begin{tcolorbox}[
  enhanced,
  colback=orange!4!white,
  colframe=orange!80!black,
  arc=2mm, boxrule=0.8pt,
  left=2mm, right=2mm, top=1.5mm, bottom=1.5mm,
  title={\footnotesize\bfseries\color{white}\faSitemap\ SƠ ĐỒ TƯ DUY: BỐN NHÁNH ĐO LƯỜNG TRẢI NGHIỆM}
]
\centering
\begin{tikzpicture}[
  box/.style={draw, rounded corners=2pt, align=center, font=\scriptsize\bfseries, inner sep=2.5pt, minimum height=0.75cm},
  arr/.style={-{Stealth[length=2mm]}, thick}
]
  \node[box, fill=purple!15, draw=purple!80!black, text width=2.6cm] (e1) at (0, 0) {ĐỘ DÀI\\$mm, cm, dm, m, km$\\Dụng cụ: Thước kẻ};
  \node[box, fill=teal!15, draw=teal!80!black, text width=2.6cm] (e2) at (3.5, 0) {KHỐI LƯỢNG\\$g, kg$\\Dụng cụ: Cân};
  \node[box, fill=blue!12, draw=blue!80!black, text width=2.6cm] (e3) at (7.0, 0) {DUNG TÍCH\\$ml, l$\\Dụng cụ: Ca đong};
  \node[box, fill=orange!15, draw=orange!80!black, text width=2.6cm] (e4) at (10.5, 0) {NHIỆT ĐỘ\\$^\circ\text{C}$\\Dụng cụ: Nhiệt kế};

  \draw[arr, purple!80!black] (e1) -- (e2);
  \draw[arr, teal!80!black] (e2) -- (e3);
  \draw[arr, blue!80!black] (e3) -- (e4);
\end{tikzpicture}
\end{tcolorbox}

\vspace{-2.5mm}
% --- KHỐI 4: VÍ DỤ & MẸO TRÁNH BẪY ---
\noindent
\begin{minipage}[t]{0.49\textwidth}\vspace{0pt}
\begin{tcolorbox}[
  enhanced,
  colback=green!4!white,
  colframe=green!60!black,
  arc=2mm, boxrule=0.8pt,
  left=2mm, right=2mm, top=1.8mm, bottom=1.8mm,
  title={\footnotesize\bfseries\color{white}\faCheckCircle\ Ví dụ thực hành chuẩn mực}
]
{\footnotesize
\textbf{Điền đơn vị đo thích hợp vào chỗ chấm:}\\
-- Bề dày của chiếc thẻ học sinh khoảng $1$\textbf{ mm}.\\
-- Bút chì dài khoảng $18$\textbf{ cm}.\\
-- Quả bí ngô nặng khoảng $3$\textbf{ kg}.\\
-- Hộp sữa tươi học đường có dung tích $180$\textbf{ ml}.\\
-- Nhiệt độ nước sôi là $100$\textbf{ $^\circ\text{C}$}.
}
\end{tcolorbox}
\end{minipage}
\hfill
\begin{minipage}[t]{0.49\textwidth}\vspace{0pt}
\begin{tcolorbox}[
  enhanced,
  colback=red!4!white,
  colframe=red!70!black,
  arc=2mm, boxrule=0.8pt,
  left=2mm, right=2mm, top=1.8mm, bottom=1.8mm,
  title={\footnotesize\bfseries\color{white}\faExclamationTriangle\ Mẹo tránh bẫy trắc nghiệm}
]
{\footnotesize
\begin{itemize}\setlength{\itemsep}{1.5pt}\setlength{\leftskip}{-2.5mm}
  \item \textbf{Bẫy chọn đơn vị vô lý so với đời sống:}\\
  Ví dụ chọn "bạn An nặng $30\text{ g}$" là hoàn toàn sai, phải là $30\text{ kg}$! Gói kẹo nặng $200\text{ g}$, không thể là $200\text{ kg}$.
  \item \textbf{Lời khuyên vàng:}\\
  Hãy liên hệ ngay với đồ vật thực tế xung quanh em để kiểm chứng tính hợp lý của con số và đơn vị trước khi chọn đáp án!
\end{itemize}
}
\end{tcolorbox}
\end{minipage}


\end{document}
